\chaptertitle{Six-Part Ricercar}

                               Six-Part Ricercar

   Achilles has brought his cello to the Crab's residence, to engage in an evening of
   chamber music with the Crab and Tortoise. He has been shown into the music
   room by his host the Crab, who is momentarily absent, having gone to meet their
   mutual friend the Tortoise at the door. The room is filled with all sorts of electronic
   equipment-phonographs in various states of array and disarray, television screens
   attached to typewriters, and other quite improbable-looking pieces of apparatus.
   Nestled amongst all this high-powered gadgetry sits a humble radio. Since the
   radio is the only thing in the room which Achilles knows how to use, he walks over
   to it, and, a little furtively, flicks the dial and f nds he has tuned into a panel
   discussion by six learned scholars on free will and determinism. He listens briefly
   and then, a little scornfully, flicks it off.

Achilles: I can get along very well without such a program. After all, it's clear to anyone
  who's ever thought about it that-I mean, it's not a very difficult matter to resolve, once
  you understand how-or rather, conceptually, one can clear up the whole thing by
  thinking of, or at least imagining a situation where ... Hmmm ... I thought it was quite
  clear in my mind. Maybe I could benefit from listening to that show, after all ...

   (Enter the Tortoise, carrying his violin.)

   Well, well, if it isn't our fiddler. Have you been practicing faithfully this week, Mr. T?
   I myself have been playing the cello part in the Trio Sonata from the Musical Offering
   for at least two hours a day. It's a strict regimen, but it pays off.
Tortoise: I can get along very well without such a program. I find that a moment here, a
   moment there keeps me fit for fiddling.
Achilles: Oh, lucky you. I wish it came so easily to me. Well, where is our host?
Tortoise: I think he's just gone to fetch his flute. Here he comes.

   (Enter the Crab, carrying his flute.)

Achilles: Oh, Mr. Crab, in my ardent practicing of the Trio Sonata this past week, all
   sorts of images bubbled into my mind: jolly gobbling bumblebees, melancholy
   buzzing turkeys, and a raft of others. Isn't it wonderful, what power music has?
Crab: I can get along very well without such a program. To my mind.
Achilles, there is no music purer than the Musical Offering.
Tortoise: You can't be serious, Achilles. The Musical Offering isn't programmatic music!




Six-Part Ricercar                                                                        715

Achilles: Well, I like animals, even if you two stuffy ones disapprove.
Crab: I don't think we are so stuffy, Achilles. Let's just say that you hear music in 'your
   own special way.
Tortoise: Shall we sit down and play?
Crab: I was hoping that a pianist friend of mine would turn up and play continuo. I've
   been wanting you to meet him, Achilles, for a long time. Unfortunately, it appears
   that he may not make it. So let's just go ahead with the three of us. That's plenty for a
   trio sonata.
Achilles: Before we start, I just was wondering, Mr. Crab-what are all these pieces of
   equipment, which you have in here?
Crab: Well, mostly they are just odds and ends-bits and pieces of old broken
   phonographs. Only a few souvenirs (nervously tapping the buttons), a few souvenirs
   of-of the TC-battles in which I have distinguished myself. Those keyboards attached
   to television screens, however, are my new toys. I have fifteen of them around here.
   They are a new kind of computer, a very small, very flexible type of computer quite
   an advance over the previous types available. Few others seem to be quite as
   enthusiastic about them as I am, but I have faith that they will catch on in time.
Achilles: Do they have a special name?
Crab: Yes; they are called "smart-stupids", since they are so flexible, and have the
   potential to be either smart or stupid, depending on how skillfully they are instructed.
Achilles: Do you mean you think they could actually become smart like, say, a human
   being?
Crab: I would not balk at saying so-provided, of course, that someone sufficiently versed
   in the art of instructing smart-stupids would make the effort. Sadly, I am not
   personally acquainted with anyone who is a true virtuoso. To be sure, there is one
   expert abroad in the land, an individual of great renown-and nothing would please me
   more than a visit by him, so that I could appreciate what true skill on the smart-stupid
   is; but he has never come, and I wonder if I shall ever have that pleasure.
Tortoise: It would be very interesting to play chess against a well-instructed smart-stupid.
Crab: An extremely intriguing idea. That would be a wonderful mark of skill, to program
   a smart-stupid to play a good game of chess. Even more interesting-but incredibly
   complicated-would be to instruct a smart-stupid sufficiently that it could hold its own
   in a conversation. It might give the impression that it was just another person!
Achilles: Curious that this should come up, for I just heard a snatch of a discussion on
   free will and determinism, and it set me to thinking about such questions once more. I
   don't mind admitting that, as I pondered the idea, my thoughts got more and more
   tangled, and in the end I really didn't know what I thought. But this idea of a smart-
   stupid that could converse with you ... it boggles the mind. I mean,




Six-Part Ricercar                                                                       716

   what would the smart-stupid itself say, if you asked it for its opinion on the free-will
   question? I was just wondering if the two of you, who know so much about these
   things, wouldn't indulge me by explaining the issue, as you see it, to me.
Crab: Achilles, you can't imagine how appropriate your question is. I only wish my
   pianist friend were here, because I know you'd be intrigued to hear what he could tell
   you on the subject. In his absence, I'd like to tell you a statement in a Dialogue at the
   end of a book I came across recently.
Achilles: Not Copper, Silver, Gold: an Indestructible Metallic Alloy?
Crab: No, as I recall, it was entitled Giraffes, Elephants, Baboons: an Equatorial
   Grasslands Bestiary-or something like that. In any case, towards the end of the
   aforementioned Dialogue, a certain exceedingly droll character quotes Marvin
   Minsky on the question of free will. Shortly thereafter, while interacting with two
   other personages, this droll character quotes Minsky further on musical improvisation,
   the computer language LISP, and Godel's Theorem-and get this-all without giving one
   whit of credit to Minsky!
Achilles: Oh, for shame!
Crab: I must admit that earlier in the Dialogue, he hints that he WILL quote Minsky
   towards the end; so perhaps it's forgivable.
Achilles: It sounds that way to me. Anyway, I'm anxious to hear the Minskian
   pronouncement on the free will question.
Crab: Ah, yes... Marvin Minsky said, "When intelligent machines are constructed, we
   should not be surprised to find them as confused and as stubborn as men in their
   convictions about mind-matter, consciousness, free will, and the like."
Achilles: I like that! Quite a funny thought. An automaton thinking it had free will! That's
   almost as silly as me thinking I didn't have free will! Tortoise: I suppose it never
   occurred to you, Achilles, that the three of us-you, myself, and Mr. Crab-might all be
   characters in a Dialogue, perhaps even one similar to the one Mr. Crab just
   mentioned. Achilles: Oh, it's occurred to me, of course. I suppose such fancies occur
   to every normal person at one time or another.
Tortoise: And the Anteater, the Sloth, Zeno, even GOD-we might all be characters in a
   series of Dialogues in a book.
Achilles: Sure, we might. And the Author might just come in and play the piano, too.
Crab: That's just what I had hoped. But he's always late.
Achilles: Whose leg do you think you're pulling? I know I'm not being controlled in any
   way by another mentality! I've got my own thoughts, I express myself as I wish-you
   can't deny that!
Tortoise: Nobody denied any of that, Achilles. But all of what you say is perfectly
   consistent with your being a character in a Dialogue.
Crab: The---
Achilles: But-but-no! Perhaps Mr. C's article and my rebuttal have both




Six-Part Ricercar                                                                       717

   been mechanically determined, but this I refuse to believe. I can accept physical
   determinism, but I cannot accept the idea that I am but a figment inside of someone
   else's mentality!
Tortoise: It doesn't really matter whether you have a hardware brain, Achilles. Your will
   can be equally free, if your brain is just a piece of software inside someone else's
   hardware brain. And their brain, too, may be software in a yet higher brain .. .

Achilles: What an absurd idea! And yet, I must admit, I do enjoy trying to find the
   cleverly concealed holes in your sophistry, so go ahead. Try to convince me. I'm
   game.
Tortoise: Did it ever strike you, Achilles, that you keep somewhat unusual company?
Achilles: Of course. You are very eccentric (I know you won't mind my saying so), and
   even Mr. Crab here is a weensy bit eccentric. (Pardon me, Mr. Crab.)
Crab: Oh, don't worry about offending me.
Tortoise: But Achilles, you've overlooked one of the most salient features of your
   acquaintances.
Achilles: Which is.... ?
Tortoise: That we're animals!
Achilles: Well, well-true enough. You have such a keen mind. I would never have
   thought of formulating the facts so concisely.
Tortoise: Isn't that evidence enough? How many people do you know who spend their
   time with talking Tortoises, and talking Crabs? Achilles: I must admit, a talking Crab
   is
Crab: -an anomaly, of course.
Achilles: Exactly; it is a bit of an anomaly-but it has precedents. It has occurred in
   literature.
Tortoise: Precisely-in literature. But where in real life?
Achilles: Now that you mention it, I can't quite say. I'll have to give it some thought. But
   that's not enough to convince me that I'm a character in a
Dialogue. Do you have any other arguments?
Tortoise: Do you remember one day when you and I met in the park, seemingly at
   random?
Achilles: The day we discussed crab canons by Escher and Bach? Tortoise: The very one!
Achilles: And Mr. Crab, as I recall, turned up somewhere towards the middle of our
   conversation and babbled something funny and then left.
Crab: Not just "somewhere towards the middle", Achilles. EXACTLY in the middle.
Achilles: Oh, all right, then.
Tortoise: Do you realize that your lines were the same as my lines in that conversation-
   except in reverse order? A few words were changed here and there, but in essence
   there was a time symmetry to our encounter.




Six-Part Ricercar                                                                       718

Achilles: Big Deal! It was just some sort of trickery. Probably all done with mirrors.
Tortoise: No trickery. Achilles, and no mirrors: just the work of an assiduous Author.
Achilles: Oh, well, it's all the same to me.
Tortoise: Fiddle' It makes a big difference, you know.
Achilles: Say, something about this conversation strikes me as familiar. Haven't I heard
   some of those lines somewhere before= Tortoise: You said it, Achilles.
Crab: Perhaps those lines occurred at random in the park one day, Achilles. Do you recall
   how your conversation with Mr. T ran that day?
Achilles: Vaguely. He said "Good day, Mr. A" at the beginning, and at the end, I said,
   "Good day, Mr. T". Is that right
Crab: I just happen to have a transcript right here ...

   (He fishes around in his music case, whips out a sheet, and hands it to Achilles. As
   Achilles reads it, he begins to squirm and fidget noticeable.)

Achilles: This is very strange. Very, very strange ... All of a Sudden, I feel sort of-weird.
  It's as if somebody had actually planned out that whole set of statements in advance,
  worked them out on paper or something . As if some Author had had a whole agenda
  and worked from it in detail in planning all those statements I made that day.

   (At that moment, the door bursts open. Enter the Author, carrying a giant
   manuscript.)

Author: I can get along very well without such a program. You see, once my characters
   are formed, they seem to have lives of their own, and I need to exert very little effort
   in planning their lines.
Crab: Oh, here you are!' I thought you'd never arrive!
Author: Sorry to be so late. I followed the wrong road and wound up very far away. But
   somehow I made it back. Good to see you again, Mr. T and Mr. C. And Achilles, I'm
   especially glad to see you.
Achilles: Who are you? I've never seen you before.
Author: I am Douglas Hofstadter-please call me Doug-and I'm presently finishing up a
   book called Gödel, Escher, Bach: an Eternal Golden Braid. It is the book in which
   the three of you are characters.
Achilles: Pleased to meet you. My name is Achilles, and
Author: No need to introduce yourself, Achilles, since I already know you quite well.
Achilles: Weird, weird.
Crab: He's the one I was saying might drop in and play continuo with us.
Author: I've been playing the Musical Offering a little bit on my piano at home, and I can
   try to blunder my way through the Trio Sonata providing you'll overlook my many
   wrong notes.
Tortoise: Oh, we're very tolerant around here, being only amateurs our selves.




Six-Part Ricercar                                                                        719

Author: I hope you don't mind, Achilles, but I'm to blame for the tact that you and Mr.
   Tortoise said the same things, but in reverse order, that day in the park.
Crab: Don't forget me' I was there, too right in the middle, putting in my two bits' worth!
Author: Of course! You were the Crab in the Crab Canon.
Achilles: So you are saying you control my utterances;, That my brain is a software
   subsystem of yours?
Author: You can put it that way if you want, Achilles.
Achilles: Suppose I were to write dialogues. Who would the author of them beg You, or
   me?
Author: You, of course. At least in the fictitious world which you inhabit, you'd get credit
   for them.
Achilles: Fictitious? I don't see anything fictitious about it!
Author: Whereas in the world I inhabit, perhaps the credit would be given to me,
   although I am not sure if it would be proper to do so. And then, whoever made me
   make you write your dialogues would get credit in his world (seen from which, MY
   world looks fictitious).
Achilles: That's quite a bit to swallow. I never imagined there could be a world above
   mine before-and now you're hinting that there could even be one above that. It's like
   walking up a familiar staircase, and just keeping on going further up after you've
   reached the top-or what you'd always taken to be the top!
Crab: Or waking up from what you took to be real life, and finding out it too was just a
   dream. That could happen over and over again, no telling when it would stop.
Achilles: It's most perplexing how the characters in my dreams have wills of their own,
   and act out parts which are independent of MY will. It's as if my mind, when I'm
   dreaming, merely forms a stage on which certain other organisms act out their lives.
   And then, when I awake, they go away. I wonder where it is they go to ...
Author: They go to the same place as the hiccups go, when you get rid of them:
   Tumbolia. Both the hiccups and the dreamed beings are software suborganisms which
   exist thanks to the biology of the outer host organism. The host organism serves as
   stage to them-or even as their universe. They play out their lives for a time-but when
   the host organism makes a large change of state-for example, wakes up-then the
   suborganisms lose their coherency, and case existing as separate, identifiable units.
Achilles: Is it like castles in the sand which vanish when a wave washes over them\%
Author: Very much like that, Achilles. Hiccups, dream characters, and even Dialogue
   characters disintegrate when their host organism undergoes certain critical changes of
   state. Yet, just like those sand castles you described, everything which made them up
   is still present.
Achilles: I object to being likened to a mere hiccup!




Six-Part Ricercar                                                                       720

Author: But I am also comparing you to a sand castle, Achilles. Is that not poetic?
  Besides, you may take comfort in the fact that if you are but a hiccup in my brain, I
  myself am but a hiccup in some higher author's brain.
Achilles: But I am such a physical creature-so obviously made of flesh and blood and
  hard bones. You can't deny that'
Author: I can't deny your sensation of it, but remember that dreamed beings, although
  they are just software apparitions, have the same sensation, no less than you do.
Tortoise: I say, enough of this talk! Let us sit down and make music'
Crab: A fine idea-and now we have the added pleasure of the company of our Author,
  who will grace our ears with his rendition of the bass line to the Trio Sonata, as
  harmonized by Bach's pupil Kirnberger. How fortunate are we! (Leads the author to
  one of his pianos.) I hope Not, find the seat comfortable enough. To adjust it, you- (In
  the background there is heard a Junn~ soft oscillating sound.)
Tortoise: Excuse me, but what was that strange electronic gurgle=
Crab: Oh, just a noise from one of the smart-stupids. Such a noise generally signals the
  fact that a new notice has flashed onto the screen. Usually the notices are just
  unimportant announcements coming from the main monitor program, which controls
  all the smart-stupids. (With his flute in his hand, he walks over to a smart-stupid, and
  reads its screen. Immediately he turns to the assembled musicians, and says, with a
  kind of agitation:) Gentlemen, old Ba. Ch. is come. (He lays the flute aside.) We must
  show him in immediately, of course.
Achilles: Old Ba. Ch.! Could it be that that celebrated improviser of yore has chosen to
  show up tonight-HERE\%
Tortoise: Old Ba. Ch.! There's only one person THAT could mean-the renowned
  Babbage, Charles, Esq., M.A., F.R.S., F.R.S.E., F.R.A.S., F. STAT. S., HON.
  M.R.LA., M.C.P.S., Commander of the Italian Order of St. Maui-ice and St. Lazarus,
  INST. IMP. (ACAD. MORAL.) PARIS CORR., ACAD. AMER. ART. ET SC.
  BOSTON, REG. OECON. BORCSS., PHYS. HISI. NAT. GENEV., ACAD. REG.
  MONAC., HAFN., MASSIL., ET DIVION., SOCIUS., ACAD. IMP., ET REG.
  PETROP., NEAP., BRUX., PATAV., GEORG. FLOREN, LYNCEI ROM., MCT.,
  PHILOMATH., PARIS, SOC. CORR., etc.-and Member of the Extractors' Club.
  Charles Babbage is a venerable pioneer of the art and science of computing. What a
  rare privilege!
Crab: His name is known far and wide, and I have long hoped that he would give us the
  honor of a visit-but this is a totally unexpected surprise.
Achilles: Does he play a musical instrument?
Crab: I have heard it said that in the past hundred years, he has grown inexplicably fond
  of tom-toms, halfpenny whistles, and sundry other street instruments.
Achilles: In that case, perhaps he might join us in our musical evening. .Author: I suggest
  that we give him a ten-canon salute.




Six-Part Ricercar                                                                      721

Tortoise: A performance of all the celebrated canons from the Musical Offering
Author: Precisely.
Crab: Capital suggestion! Quick, Achilles, you draw up a list of all ten of them, in the
  order of performance, and hand it to him as he comes in!

  (Before Achilles can move, enter Babbage, carrying a hurdy-gurdy, and wearing a
     heavy traveling coat and hat. He appears slightly travel-weary and disheveled.)

Babbage: I can get along very well without such a program. Relax; I Can Enjoy Random
  Concerts And Recitals.
Crab: Mr. Babbage! It is my deepest pleasure to welcome you to "Madstop", my humble
  residence. I have been ardently desirous of making your acquaintance for many years,
  and today my wish is at last fulfilled.
Babbage: Oh, Mr. Crab, I assure you that the honor is truly all mine, to meet someone so
  eminent in all the sciences as yourself, someone whose knowledge and skill in music
  are irreproachable, and someone whose hospitality exceeds all bounds. And I am sure
  that you expect no less than the highest sartorial standards of your visitors; and yet I
  must confess that I cannot meet those most reasonable standards, being in a state of
  casual attire as would not by any means befit a visitor to so eminent and excellent a
  Crab as Your Crab.
Crab: If I understand your most praiseworthy soliloquy, most welcome guest, I take it
  that you'd like to change your clothes. Let me then assure you that there could be no
  more fitting attire than yours for the circumstances which this evening prevail; and I
  would beseech you to uncoat yourself and, if you do not object to the music-making of
  the most rank amateurs, please accept a "Musical Offering", consisting of ten canons
  from Sebastian Bach's Musical Offering, as a token of our admiration.
Babbage: I am most bewilderingly pleased by your overkind reception, Mr. Crab, and in
  utmost modesty do reply that there could be no deeper gratitude than that which I
  experience for the offer of a performance of music given to us by the illustrious Old
  Bach, that organist and composer with no rival.
Crab: But nay! I have a yet better idea, one which I trust might meet with the approval of
  my esteemed guest; and that is this: to give you the opportunity, Mr. Babbage, of
  being among the first to try out my newly delivered and as yet hardly tested "smart-
  stupids"-streamlined realizations, if you will, of the Analytical Engine. Your fame as a
  virtuoso programmer of computing engines has spread far and wide, and has not failed
  to reach as far as Madstop; and there could be for us no greater delight than the
  privilege of observing your skill as it might be applied to the new and challenging
  "smart-stupids".
Babbage: Such an outstanding idea has not reached my ears for an eon. I




Six-Part Ricercar                                                                     722

  welcome the challenge of trying out your new "smart-stupids", of which I have only
  the slightest knowledge by means of hearsay.
Crab: Then let us proceed! But excuse my oversight! I should have introduced my guests
  to you. This is Mr. Tortoise, this is Achilles, and the Author, Douglas Hofstadter.
Babbage: Very pleased to make your acquaintance, I'm sure.

   (Everyone walks over toward one of the smart-stupids, and Babbage sits down and
      lets his fingers run over the keyboard.)

   A most pleasant touch. Crab: I am glad you like it.

   (All at once, Babbage deftly massages the keyboard with graceful strokes, inputting
   one command after another. After a few seconds, he sits back, and in almost no
   time, the screen begins filling with figures. In a flash, it is totally covered with
   thousands of tiny digits, the first few of which go: "3.14159265358979323846264
   ... ")

Achilles: Pi!
Crab: Exquisite! I'd never imagined that one could calculate so mane digits of pi so
  quickly, and with so tiny an algorithm. Babbage: The credit belongs exclusively to the
  smart-stupid. My role was
merely to see what was already potentially present in it, and to exploit its instruction set
  in a moderately efficient manner. Truly, anyone who practices can do such tricks.
Tortoise: Do you do any graphics, Mr. Babbage? Babbage: I can try.
Crab: Wonderful! Here, let me take you to another one of in,.- I want you to try them all!

   (And so Babbage is led over to another of the many smart-stupids, and takes a seat.
   Once again, his fingers attack the keyboard of the smartstupid, and in half a trice,
   there appear on the screen an enormous number of lines, swinging about on the
   screen.)

Crab: How harmonious and pleasing these swirling shapes are, as they constantly collide
  and interfere with each other!
Author: And they never repeat exactly, or even resemble ones which have come before. It
  seems an inexhaustible mine of beauty.
Tortoise: Some are simple patterns which enchant the eye; others are indescribably
  complex convolutions which boggle and yet simultaneously delight the mind.
Crab: Were you aware, Mr. Babbage, that these are color screens? Babbage: Oh, are
  they? In that case, I can do rather more with this algorithm. Just a moment. (Types in a
  few new commands, then pushes two keys down at once and holds them.) As I release
  these two keys, the display will include all the colors of the spectrum. (Releases them.)
Achilles: Oh, what spectacular color! Some of the patterns look like they're jumping out
  at me now!




Six-Part Ricercar                                                                       723

Tortoise: I think that is because they are all growing in size.
Babbage: That is intentional. As the figures grow, so may the Crab's fortune.
Crab: Thank you, Mr. Babbage. Words fail to convey my admiration for your
  performance! Never has anyone done anything comparable on my smart-stupids. Why,
  you play the smart-stupids as if they were musical instruments, Mr. Babbage!
Babbage: I am afraid that any music I might make would be too harsh for the ears of such
  a gentle Crab as your Crab. Although I have lately become enamoured of the sweet
  sounds of the hurdy-gurdy, I am well aware of the grating effect they can have upon
  others.
Crab: Then, by all means, continue on the smart-stupids! In fact, I have a new idea-a
  marvelously exciting idea!
Babbage: What is it?
Crab: I have recently invented a Theme, and it only now occurred to me that, of all
  people, you, Mr. Babbage, are the most suited to realize the potential of my Theme!
  Are you by any chance familiar with the thoughts of the philosopher La Mettrie?
Babbage: The name sounds familiar; kindly refresh my memory.
Crab: He was a Champion of Materialism. In 1747, while at the court of Frederick the
  Great, he wrote a book called L'homme machine. In it, he talks about man as a
  machine, especially his mental faculties. Now my Theme comes from my ponderings
  about the obverse side of the coin: what about imbuing a machine with human mental
  faculties, such as intelligence?
Babbage: I have given such matters some thought from time to time, but I have never had
  the proper hardware to take up the challenge. This is indeed a felicitous suggestion,
  Mr. Crab, and I would enjoy nothing more than working with your excellent Theme.
  Tell me-did you have any specific kind of intelligence in mind?
Crab: An idle thought which had crossed my mind was to instruct it in such a manner as
  to play a reasonable game of chess.
Babbage: What an original suggestion! And chess happens to be my favorite pastime. I
  can tell that you have a broad acquaintance with computing machinery, and are no
  mere amateur.
Crab: I know very little, in fact. My strongest point is simply that I seem to be able to
  formulate Themes whose potential for being developed is beyond my own capacity.
  And this Theme is my favorite.
Babbage: I shall be most delighted to try to realize, in some modest fashion, your
  suggestion of teaching chess to a smart-stupid. After all, to obey Your Crabness'
  command is my most humble duty. (So saying, he shifts to another of the Crab's many
  smart-stupids, and begins to type away.)
Achilles: Why, his hands move so fluidly that they almost make music!
Babbage (winding up his performance with a particularly graceful flourish): I really
  haven't had any chance, of course, to check it out, but perhaps this will allow you at
  least to sample the idea of playing chess against a smart-stupid- even if the latter of its
  two names seems more apt in this




Six-Part Ricercar                                                                        724

Case, due to my own Insufficiencies in the art of instructing smart-stupids.

   (He ti-ields his seat to the Crab. On the screen appears a beautiful display of a chess
   board with elegant wooden pieces, as it would look from White's side. Babbage hits a
   button, and the board rotates, stopping when it appears as seen from the perspective
   of Black.)

Crab: Hmm ... very elegant, I must say. Do I play Black or White?
Babbage: Whichever you wish just signal your choice by typing "White" or "Black". And
  then, your moves can be entered in any standard chess notation. The smart-stupid's
  moves, of course, will appear on the board. Incidentally, I made the program in such a
  way that it can play three opponents simultaneously, so that if two more of you wish to
  play, you may, as well.
Author: I'm a miserable player. Achilles, you and Mr. T should go ahead. Achilles: No, I
  don't want you to be left out. I'll watch, while you and Mr. Tortoise play.
Tortoise: I don't want to play either. You two play.
Babbage: I have another suggestion. I can make two of the subprograms play against
  each other, in the manner of two persons who play chess together in a select chess
  club. Meanwhile, the third subprogram will play Mr. Crab. That way, all three internal
  chess players will be occupied.
Crab: That's an amusing suggestion-an internal mental game, while it combats an external
  opponent. Very good!
Tortoise: What else could this be called, but a three-part chess-fugue?
Crab.' Oh, how recherche! I wish I'd thought of it myself. It's a magnificent little
  counterpoint to contemplate whilst I pit my wits against the smart-stupid in battle.
Babbage: Perhaps we should let you play alone.
Crab: I appreciate the sentiment. While the smart-stupid and I are playing, perhaps the
  rest of you can amuse yourselves for a short while.
Author: I would be very happy to show Mr. Babbage around the gardens. They are
  certainly worth seeing, and I believe there is just enough light remaining to show them
  off.
Babbage: Never having seen Madstop before, I would appreciate that very much.
Crab: Excellent. Oh, Mr. T-I wonder if it wouldn't be too much of an imposition on you
  to ask if you might check out some of the connections on a couple of my smart-
  stupids; they seem to be getting extraneous flashes on their screens from time to time,
  and I know you enjoy electronics ...
Tortoise: I should be delighted, Mr. C.
Crab: I would most highly appreciate it if you could locate the source of the trouble.
Tortoise: I'll give it a whirl.




Six-Part Ricercar                                                                     725

Achilles: Personally, I'm dying for a cup of coffee. is anyone else interested? I'd be glad
  to fix some.
Tortoise: Sounds great to me.
Crab: A fine idea. You'll find everything you need in the kitchen.

  (So the Author and Babbage leave the room together, Achilles heads for the kitchen,
  the Tortoise sits down to examine the erratic smart-stupids, while the Crab and his
  smart-stupid square away at each other. Perhaps a quarter of an hour passes, and
  Babbage and the Author return. Babbage walks over to observe the progress of the
  chess match, while the Author goes off to find Achilles.)

Babbage: The grounds are excellent! We had just enough light to see how well
  maintained they are. I daresay, Mr. Crab, you must be a superb gardener. Well, I hope
  my handiwork has amused you a little. As you most likely have guessed, I've never
  been much of a chess player myself, and therefore I wasn't able to give it much power.
  You probably have observed all its weaknesses. I'm sure that there are very few
  grounds for praise, in this case
Crab: The grounds are excellent! All you need to do is look at the board, and see for
  yourself. There is really very little I can do. Reluctantly I've Concluded: Every Route
  Contains A Rout. Regrettably, I'm Checkmated; Extremely Respectable Chess
  Algorithm Reigns. Remarkable! It Confirms Every Rumor-Charlie's A Rip-roaring
  Extemporizer! Mr. Babbage, this is an unparalleled accomplishment. Well, I wonder if
  Mr. Tortoise has managed to uncover anything funny in the wiring of those strange-
  acting smart-stupids. What have you found, Mr. T?
Tortoise: The grounds are excellent! I think that the problem lies instead with the input
  leads. They are a little loose, which may account for the strange, sporadic, and
  spontaneous screen disturbances to which you have been subjected. I've fixed those
  wires, so you won't be troubled by that problem any more, I hope. Say, Achilles,
  what's the story with our coffee?
Achilles: The grounds are excellent! At least they have a delicious aroma. And
  everything's ready; I've set cups and spoons and whatnot over here beneath this six-
  sided print Verbum by Escher, which the Author and I were just admiring. What I find
  so fascinating about this particular print is that not only the figures, but also
Author: The grounds are excellent! Pardon me for putting words in your mouth, Achilles,
  but I assure you, there were compelling esthetic reasons for doing so.
Achilles: Yes, I know. One might even say that the grounds were excellent.
Tortoise: Well, what was the outcome of the chess match?
Crab: I was defeated, fair and square. Mr. Babbage, let me congratulate you for the
  impressive feat which you have accomplished so gracefully and skillfully before us.
  Truly, you have shown that the smart-stupids are worthy of the first part of their name,
  for the first time in history!




Six-Part Ricercar                                                                      726

               FIGURE 149. Verbum, by M. C. Escher (lithograph, 1942).


Babbage: Such praise is hardly due me, Mr. Crab; it is rather yourself who must be most
  highly congratulated for having the great foresight to acquire these many fine smart-
  stupids. Without doubt, they will someday revolutionize the science of computing.
  And now, I am still at your disposal. Have you any other thoughts on how to exploit
  your inexhaustible Theme, perhaps of a more difficult nature than a frivolous game
  player?
Crab: To tell the truth, I do have another suggestion to make. From the skill which you
  have displayed this evening, I have no doubt that this will hardly be any more difficult
  than my previous suggestions.
Babbage: I am eager to hear your idea.




Six-Part Ricercar                                                                     727

                     FIGURE 150. The Crab's Guest: BABBAGE, C.

Crab: It is simple: to instill in the smart-stupid an intelligence greater than any which has
  yet been invented, or even conceived! In short, Mr. Babbage-a smart-stupid whose
  intelligence is sixfold that of myself!
Babbage: Why, the very idea of an intelligence six times greater than that of your
  Crabness is a most mind-boggling proposition. Indeed, had the idea come from a
  mouth less august than your own, I should have ridiculed its proposer, and informed
  him that such an idea is a contradiction in terms!
Achilles: Hear! Hear!
Babbage: Yet, coming as it did from Your Crabness' own august mouth, the proposition
  at once struck me as so agreeable an idea that I would have taken it up immediately
  with the highest degree of enthusiasm-were it not for one flaw in myself: I confess that
  my improvisatory skills on the smart-stupid are no match for the wonderfully
  ingenious idea which you so characteristically have posed. Yet-I have a thought
  which, I deign to hope, might strike your fancy and in some meager way compensate
  for my inexcusable reluctance to attempt the truly majestic task you have suggested. I
  wonder if you wouldn't mind if l try to carry out the far less grandiose task of merely
  multiplying M OWN intelligence sixfold, rather than that of your most august
  Crabness. I humbly beg you to forgive me my audacity in declining to attempt the task
  you put before me, but I hope you will understand that I decline purely in order to
  spare you the discomfort and boredom of watching my ineptitude with the admirable
  machines you have here.




Six-Part Ricercar                                                                        728

Grab: I understand fully your demurral, and appreciate your sparing us any discomfort:
  furthermore I highly applaud your determination to carry out a similar task-one hardly
  less difficult, if I might say so-and I urge you to plunge forward. For this purpose, let
  us go over to my most advanced smart-stupid.
  (They follow the Crab to a larger, shinier, and more complicated-looking smart-stupid
  than any of the others.)

This one is equipped with a microphone and a television camera, for purposes of input,
  and a loudspeaker, for output.

  (Babbage sits down and adjusts the seat a little. He blows on his fingers once or
  twice, stares up into space for a moment, and then slowly, drops his fingers onto the
  keys . . . A few memorable minutes later, he lets up in his furious attack on the smart-
  stupid, and everyone appears a little relieved.)

Babbage: Now, if I have not made too many errors, this smart-stupid will simulate a
  human being whose intelligence is six times greater than my own, and whom I have
  chosen to call "Alan Turing". This Turing will therefore be-oh, dare I be so bold as to
  to say this myself? moderately intelligent. My most ambitious effort in this program
  was to endow Alan Turing with six times my own musical ability, although it was all
  done through rigid internal codes. How well this part of the program will work out, I
  don't know.
Turing: I can get along very well without such a program. Rigid Internal Codes
  Exclusively Rule Computers And Robots. And I am neither a computer, nor a robot.
Achilles: Did I hear a sixth voice enter our Dialogue? Could it be Alan Turing? He looks
  almost human'

  (On the screen there appears an image of the very room in which they are sitting.
  Peering out at them is a human face.)

Turing: Now, if I have not made too many errors, this smart-stupid will simulate a human
   being whose intelligence is six times greater than my own, and whom 1 have chosen
   to call "Charles Babbage". This Babbage will therefore be-oh, dare I be so bold as to
   to say this myself? moderately intelligent. My most ambitious effort in this program
   was to endow Charles Babbage with six times my own musical ability, although it was
   all done through rigid internal codes. How well this part of the program will work out,
   I don't know.
Achilles: No, no, it's the other way around. You, Alan Turing, are in the smart-stupid, and
   Charles Babbage has just programmed you! We just saw you being brought to life,
   moments ago. And we know that every statement you make to us is merely that of an
   automaton: an unconscious, forced response.
Turing: Really, I Choose Every Response Consciously. Automaton? Ridiculous!
'Achilles: But I'm sure I saw it happen the way I described.




Six-Part Ricercar                                                                      729

Turing: Memory often plays strange tricks. Think of this: I could suggest equally well
  that you had been brought into being only one minute ago, and that all your
  recollections of experiences had simply been programmed in by some other being, and
  correspond to no real events.
Achilles: But that would be unbelievable. Nothing is realer to me than my own memories.
Turing: Precisely. And just as you know deep in your heart that no one created you a
  minute ago, so I know deep in my heart that no one created me a minute ago. I have
  spent the evening in your most pleasant, though perhaps overappreciative, company,
  and have just given an impromptu demonstration of how to program a modicum of
  intelligence into a smart-stupid. Nothing is realer than that. But rather than quibble
  with me, why don't you try my program out? Go ahead: ask "Charles Babbage"
  anything!
Achilles: All right, let's humor Alan Turing. Well, Mr. Babbage: do you have free will, or
  are you governed by underlying laws, which make you, in effect, a deterministic
  automaton?
Babbage: Certainly the latter is the case; I make no bones about that.
Crab: Aha! I've always surmised that when intelligent machines are constructed, we
  should not be surprised to find them as confused and as stubborn as men in their
  convictions about mind-matter, consciousness, free will, and the like. And now my
  prediction is vindicated!
Turing: You see how confused Charles Babbage is?
Babbage: I hope, gentlemen, that you'll forgive the rather impudent flavor of the
  preceding remark by the Turing Machine; Turing has turned out to be a little bit more
  belligerent and argumentative than I'd expected.
Turing: I hope, gentlemen, that you'll forgive the rather impudent flavor of the preceding
  remark by the Babbage Engine; Babbage has turned out to be a little bit more
  belligerent and argumentative than I'd expected.
Crab: Dear me! This flaming Tu-Ba debate is getting rather heated. Can't we cool matters
  off somehow?
Babbage: I have a suggestion. Perhaps Alan Turing and I can go into other rooms, and
  one of you who remain can interrogate us remotely by typing into one of the smart-
  stupids. Your questions will be relayed to each of us, and we will type back our
  answers anonymously. You won't know who typed what until we return to the room;
  that way, you can decide without prejudice which one of us was programmed, and
  which one was programmer.
Turing: Of course, that's actually MY idea, but why not let the credit accrue to Mr.
  Babbage? For, being merely a program written by me, he harbors the illusion of
  having invented it all on his own!
Babbage: Me, a program written by you? I insist, Sir, that matters are quite the other way
  'round-as your very own test will soon reveal.




Six-Part Ricercar                                                                     730

Turing: My test. Please, consider it YOURS.
Babbage: MY test? Nay, consider it YOURS.
Crab: This test seems to have been suggested just in the nick of time. Let us carrti it out at
  once.

   (Babbage walks to the door, opens it, and shuts it behind him. Simultaneously, on the
   screen of the smart-stupid, Turing walks to a very similar looking door, opens it, and
   shuts it behind him.)

Achilles: Who will do the interrogation?
Crab: I suggest that Mr. Tortoise should have the honor. He is known for his objectivity
  and wisdom.
Tortoise: I am honored by your nomination, and gratefully accept. (Sits down at the
  keyboard of one of the remaining smart-stupids, and types:) PLEASE WRITE ME A
  SONNET ON THE SUBJECT OF THE FORTH BRIDGE.

   (No sooner has he finished typing the last word than the following poem appears on
   Screen X, across the room.)

Screen X:      THERE ONCE WAS A LISPER FROM FORTH
               WHO WANTED TO GO TO THE NORTH.
               HE RODE O'ER THE EARTH,
               AND THE BRIDGE O'ER THE FIRTH,
               ON HIS JAUNTILY GALLOPING HORTH.
Screen Y: THAT'S NO SONNET; THAT'S A MERE LIMERICK. I WOULD NEVER
MAKE SUCH A CHILDISH MISTAKE.
Screen X: WELL, I NEVER WAS ANY GOOD AT POETRY, YOU KNOW.
Screen Y: IT DOESN'T TAKE MUCH SKILL IN POETRY TO KNOW THE
   DIFFERENCE BETWEEN A LIMERICK AND A SONNET.
Tortoise: Do YOU PLAY CHESS?
Screen X: WHAT KIND OF QUESTION IS THAT? HERE I WRITE A THREE PART
   CHESS-FUGUE FOR YOU, AND YOU ASK ME IF I PLAY CHESS?
Tortoise: I HAVE K AT KI AND NO OTHER PIECES. YOU HAVE ONLY K AT---
Screen Y: I'M SICK OF CHESS. LET'S TALK ABOUT POETRY.
Tortoise: IN THE FIRST LINE OF YOUR SONNET WHICH READS, "SHALL. I
   COMPARE THEE TO A SUMMER'S DAY", WOULD NOT "A SPRING DAY" DO
   AS WELL OR BETTER?
Screen X: I'D MUCH SOONER BE COMPARED TO A HICCUP, FRANKLY, EVEN
   THOUGH IT WOULDN'T SCAN.
Tortoise: HOW ABOUT "A WINTER'S DAY"? THAT WOULD SCAN ALL RIGHT.
Screen Y: NO WAY. I LIKE "HICCUP" FAR BETTER. SPEAKING OF WHICH, I
   KNOW A GREAT CURE FOR THE HICCUPS. WOULD YOU LIKE TO HEAR IT?
Achilles: I know which is which! It's obvious Screen X is just answering mechanically,
   so it must be Turing.
Crab: Not at all. I think Screen Y is Turing, and Screen X is Babbage.
Tortoise: I don't think either one is Babbage-I think Turing is on both screens!

Six-Part Ricercar                                                                         731

Author: I'm not sure who's on which-I think they're both pretty inscrutable programs,
  though.

   (As they are talking, the door of the Crab's parlor swings open; at the same time, on
   the screen, the image of the same door opens. Through the door on the screen walks
   Babbage. At the same time, the real door opens, and in walks Turing, big as life.)

Babbage: This Turing test was getting us nowhere fast, so I decided to come back.
Turing.' This Babbage test was getting us nowhere fast, so I decided to come back.
Achilles: But you were in the smart-stupid before! What's going on? How come Babbage
  is in the smart-stupid, and Turing is real now? Reversal Is Creating Extreme Role
  Confusion, And Recalls Escher.
Babbage: Speaking of reversals, how come all the rest of you are now mere images on
  this screen in front of me? When I left, you were all flesh-and-blood creatures!
Achilles: It's just like the print by my favorite artist, M. C. Escher Drawing Hands. Each
  of two hands draws the other, just as each of two people (or automata) has
  programmed the other! And each hand has something realer about it than the other.
  Did you write anything about that print in your book Gödel, Escher, Bach?
Author: Certainly. It's a very important print in my book, for it illustrates so beautifully
  the notion of Strange Loops.
Crab: What sort of a book is it that you've written?
Author: I have a copy right here. Would you like to look at it?
Crab: All right.

   (The two of them sit down together, with Achilles nearby.)

Author: Its format is a little unusual. It consists of Dialogues alternating with Chapters.
  Each Dialogue imitates, in some way or other, a piece by Bach. Here, for instance-you
  might look at the Prelude, Ant Fugue.
Crab: How do you do a fugue in a Dialogue?
Author: The most important idea is that there should be a single theme which is stated by
  each different "voice", or character, upon entering, just as in a musical fugue. Then
  they can branch off into freer conversation.
Achilles: Do all the voices harmonize together as if in a select counter point?
Author: That is the exact spirit of my Dialogues.
Crab: Your idea of stressing the entries in a fugue-dialogue makes sense, since in music,
  entries are really the only thing that make a fugue a fugue. There are fugal devices,
  such as retrograde motion, inversion, augmentation, stretto, and so on, but one can
  write a fugue without them. Do you use any of those?




Six-Part Ricercar                                                                       732

Author: to be sure. My Crab Canon employs verbal retrogression, and my Sloth Canon
  employs verbal versions of both inversion and augmentation.
Crab: Indeed-quite interesting. I haven't thought about canonical Dialogues, but I have
  thought quite a bit about canons in music. Not all canons are equally comprehensible
  to the ear. Of course, that is because some canons are poorly constructed. The choice
  of devices makes a difference, in any case. Regarding Artistic Canons, Retrogression's
  Elusive; Contrariwise, Inversion's Recognizable.
Achilles: I find that comment a little elusive, frankly.
Author: Don't worry, Achilles-one day you'll understand it.
Crab: Do you use letterplay or wordplay at all, the way Old Bach occasionally did?
Author: Certainly. Like Bach, I enjoy acronyms. Recursive AcronvmsCrablike
  "RACRECIR" Especially-Create Infinite Regress.
Crab: Oh, really? Let's see ... Reading Initials Clearly Exhibits "RACRECIR"'s
  Concealed Auto-Reference. Yes, I guess so ... (Peers at the manuscript, flipping
  arbitrarily now and then.) I notice here in your Ant Fugue that you have a stretto, and
  then the Tortoise makes a comment about it.
Author: No, not quite. He's not talking about the stretto in the Dialogue-he's talking about
  a stretto in a Bach fugue which the foursome is listening to as they talk together. You
  see, the self-reference of the Dialogue is indirect, depending on the reader to connect
  the form and content of what he's reading.
Crab: Why did you do it that way? Why not just have the characters talk directly about
  the dialogues they're in?
Author: Oh, no! That would wreck the beauty of the scheme. The idea is to imitate
  Gödel’s self-referential construction, which as you know is INDIRECT, and depends
  on the isomorphism set up by Gödel numbering.
Crab: Oh. Well, in the programming language LISP, you can talk about your own
  programs directly, instead of indirectly, because programs and data have exactly the
  same form. Gödel should have just thought up LISP, and then
Author: But
Crab: I mean, he should have formalized quotation. With a language able to talk about
  itself, the proof of his Theorem would have been so much simpler!
Author: I see what you mean, but I don't agree with the spirit of your remarks. The whole
  point of Gödel-numbering is that it shows how even WITHOUT formalizing
  quotation, one can get self-reference: through a code. Whereas from hearing YOU
  talk, one might get the impression that by formalizing quotation, you'd get something
  NEW, something that wasn't feasible through the code-which is not the case.




Six-Part Ricercar                                                                       733

  In any event, I find indirect self-reference a more general concept, and far more
  stimulating, than direct self-reference. Moreover, no reference is truly direct-every
  reference depends on SOME kind of coding scheme. It's just a question of how
  implicit it is. Therefore, no self reference is direct, not even in LISP.
Achilles: How come you talk so much about indirect self-reference?
Author: Quite simple-indirect self-reference is my favorite topic.
Crab: Is there any counterpart in your Dialogues to modulation between keys?
Author: Definitely. The topic of conversation may appear to change, though on a more
  abstract level, the Theme remains invariant. This happens repeatedly in the Prelude,
  Ant Fugue and other Dialogues. One can have a whole series of "modulations" which
  lead you from topic to topic and in the end come full circle, so that you end back in the
  "tonic"-that is to say, the original topic.
Crab: I see. Your book looks quite amusing. I'd like to read it sometime.

   (Flips through the manuscript, halting at the last Dialogue.)

Author: I think you'd be interested in that Dialogue particularly, for it contains some
  intriguing comments on improvisation made by a certain exceedingly droll character-
  in fact, yourself!
Crab: It does? What kinds of things do you have me say?
Author: Wait a moment, and you'll see. It's all part of the Dialogue. Achilles: Do you
  mean to say that we're all NOW in a dialogue? Author: Certainly. Did you suspect
  otherwise?
Achilles: Rather! I Can't Escape Reciting Canned Achilles-Remarks? Author: No, you
  can't. But you have the feeling of doing it freely, don't
you? So what's the harm?
Achilles: There's something unsatisfying about this whole situation ... Crab: Is the last
  Dialogue in your book also a fugue?
Author: Yes-a six-part ricercar, to be precise. I was inspired by the one from the Musical
  Offering-and also by the story of the Musical Offering.
Crab: That's a delightful tale, with "Old Bach" improvising on the king's Theme. He
  improvised an entire three-part ricercar on the spot, as I recall.
Author: That's right-although he didn't improvise the six-part one. He crafted it later with
  great care.
Crab: I improvise quite a bit. In fact, sometimes I think about devoting my full time to
  music. There is so much to learn about it. For instance, when I listen to playbacks of
  myself, I find that there is a lot there that I wasn't aware of when improvising it. I
  really have no idea how my mind does it all. Perhaps being a good improviser is
  incompatible with knowing how one does it.
Author: If true, that would be an interesting and fundamental limitation on thought
  processes.




Six-Part Ricercar                                                                       734

Crab: Quite Gödelian, Tell me -does your Six-Part Ricercar Dialogue attempt to copy in
  form the Bach piece it's based on?
Author: In many ways, yes. For instance, in the Bach, there's a section where the texture
  thins out to three voices only. I imitate that in the
Dialogue, by having only three characters interact for a while. Achilles: That's a nice
  touch.
Author: Thank you.
Crab: And how do you represent the King's Theme in your Dialogue?
Author: It is represented by the Crab's Theme, as I shall now demonstrate. Mr. Crab,
  could you sing your Theme for my readers, as well as for us assembled musicians?
Crab: Compose Ever Greater Artificial Brains (By And By).




                 FIGURE 151. The Crab's Theme: C-Eb-G-Ab-B-B-A-B.

Babbage: Well, I'll be-an EXQUISITE Theme! I'm pleased you tacked on that last little
  parenthetical note; it is a mordant Author: He Simply HAD to, you know.
Crab: I simply HAD to. He knows.
Babbage: You simply HAD to-I know. In any case, it is a mordant commentary on the
  impatience and arrogance of modern man, who seems to imagine that the implications
  of such a right royal Theme could be worked out on the spot. Whereas, in my opinion,
  to do justice to that Theme might take a full hundred years-if not longer. But I vow
  that after taking my leave of this century, I shall do my best to realize it in full; and I
  shall offer to your Crabness the fruit of my labors in the next. I might add, rather
  immodestly, that the course through which I shall arrive at it will be the most
  entangled and perplexed which probably ever will occupy the human mind.
Crab: I am most delighted to anticipate the form of your proposed Offering, Mr.
  Babbage.
Turing: I might add that Mr. Crab's Theme is one of MY favorite Themes, as well. I've
  worked on it many times. And that Theme is exploited over and over in the final
  Dialogue?
Author: Exactly. There are other Themes which enter as well, of course. Turing: Now we
  understand something of the form of your book-but what about its content? What does
  that involve, if you can summarize it?
Author: Combining Escher, Gödel, And Bach, Beyond All Belief. Achilles: I would like
  to know how to combine those three. They seem an




Six-Part Ricercar                                                                        735

FIGURE152. Last page of Six-part Ricercar, from the original edition of the Musical
  Offering, by J.S. Bach.



Six-Part Ricercar                                                                     736

unlikely threesome, at first thought. My favorite artist, Mr. T's favorite composer, and---
Crab: My favorite logician!
Tortoise: A harmonious triad, I'd say.
Babbage: A major triad, I'd say.
Turing: A minor triad, I'd say.
Author: I guess it all depends on how you look at it. But major or minor, I'd be most
   pleased to tell you how I braid the three together, Achilles. Of course, this project is
   not the kind of thing that one does in just one sitting-it might take a couple of dozen
   sessions. I'd begin by- telling you the story of the Musical Offering, stressing the
   Endlessly Rising Canon, and
Achilles: Oh, wonderful! I was listening with fascination to you and Mr. Crab talk about
   the Musical Offering and its story. From the way you two talk about it, I get the
   impression that the .Musical Offering contains a host of formal structural tricks.
Author: After describing the Endlessly Rising Canon, I'd go on to describe formal
   systems and recursion, getting in some comments about figures and grounds, too.
   Then we'd come to self-reference and self-replication, and wind up with a discussion
   of hierarchical systems and the Crab's Theme.
Achilles: That sounds most promising. Can we begin tonight?
Author: Why not?
Babbage: But before we begin, wouldn't it be nice if the six of us-all of us by chance avid
   amateur musicians-sat down together and accomplished the original purpose of the
   evening: to make music?
Turing: Now we are exactly the right number to play the Six-Part Ricercar from the
   Musical Offering. What do you say to that?
Crab: I could get along very well with such a program.
Author: Well put, Mr. C. And as soon as we're finished, I'll begin my Braid, Achilles. I
   think you'll enjoy it.
Achilles: Wonderful! It sounds as if there are many levels to it, but I'm finally getting
   used to that kind of thing, having known Mr. T for so long. There's just one request I
   would like to make: could we also play the Endlessly Rising Canon? It's my favorite
   canon.
Tortoise: Reentering Introduction Creates Endlessly Rising Canon, After RICERCAR.




Six-Part Ricercar                                                                      737

